\einheitenkopf{Einheit 2 von 5 · Säulen-, Balken- und Liniendiagramme}
\erg{14}{a) 10 \quad b) 50 \quad c) 5 \quad d) $0{,}01$ \quad e) 50 Mio.}
\erg{15}{a) ja \quad b) nein, Start 40 \quad c) ja \quad d) nein, Start $3{,}10$ \quad e) nein, Start 820}
\erg{16}{a) 20 \quad b) 60 \quad c) 40 \quad d) 80 \quad e) 55 \quad f) 65 \quad g) 3 Tage (Do, Sa, So; Dienstag mit genau 60 zählt nicht mit) \quad h) meiste: Donnerstag (80), wenigste: Montag (20) \quad i) $75-40=35$ \quad j) 24\,000 \quad k) 46\,000 \quad l) $46\,000-30\,000=16\,000$}
\erg{17}{a) 24 \quad b) 14 \quad c) Blau \quad d) $4\,^\circ$C \quad e) Januar \quad f) zwischen Februar und März \quad g) Die Temperatur steigt von März bis Juni durchgehend an, von $4\,^\circ$C auf $17\,^\circ$C.}
\erg{18}{a) Säule bis 15 \quad b) bis 25 \quad c) bis 10 \quad d) bis 30 \quad e) bis 20 \quad f) Die Achse müsste mindestens bis 80 reichen, zum Beispiel in Schritten von 20.}
\erg{19}{a) $120 : 4 = 30$ \quad b) $6 \cdot 30 = 180$ \quad c) $3 \cdot 30 = 90$ \quad d) Säule 7 Kästchen hoch ($210 : 30 = 7$) \quad e) $480 : 6 = 80$ Mio. je Kästchen; die zweite Säule (4 Kästchen) ist Hindi, die dritte (3 Kästchen) Portugiesisch; Arabisch: $400 : 80 = 5$ Kästchen}
\erg{20}{a) Jonas hat die Beschriftung „in Tausend“ übersehen. Richtig: $18 \cdot 1000 = 18\,000$ Zuschauer. \quad b) 25\,000}
\erg{21}{a) Ein Kreisdiagramm zeigt Anteile an einem Ganzen. Die Geburtenzahlen sind aber Werte zu zehn verschiedenen Zeitpunkten; Säulen oder eine Linie zeigen, wie sich der Wert über die Jahre verändert. \quad b) Nein. Vier Hilfslinien teilen den Abschnitt von 0 bis 50 in fünf gleiche Teile, eine Linie ist also 10 wert.}
